\documentclass[book.tex]{subfiles}
\begin{document}

\chapter{Long Range Atomic Contacts}


\label{chap:long_range_atomic}
\noindent\rule{11cm}{0.4pt} \\
\textbf{Prerequisites:} \autoref{chap:graphconvs}, \autoref{chap:equivariant_networks}, \autoref{chap:molecular_dynamics},  \autoref{chap:dft} \\
\textbf{Difficulty Level:} ** \\
\noindent\rule{11cm}{0.4pt}
\vspace{5mm}

In  previous chapters, we have learned about some mathematical representations of molecular structure that respects the relevant invariances of the underlying system. However, we made a key approximation by introducing cutoffs into the system that assume that all interactions happen within a certain distance. This approximation limits the accuracy of our models since longer range interactions do play a world in real molecular systems.

In this chapter, we introduce long-range features that incorporate a form of Coulomb-interactions into computed descriptors \cite{grisafi2018symmetry}, \cite{grisafi2018transferable}, \cite{grisafi2019incorporating}. This representation represents a molecules as an $O(3)$ equivariant feature vector. These longer range interactions will allow us to represent more accurately long-range physical effects such as polarization and dispersion interactions.
Representing long range contacts is challenging for atomic systems in force fields. 

%\textcolor{red}{TODO: The summary in the comments discusses a covariant representation. Do we want to preserve that as well?}

%"""Most transferable machine-learning models for molecules and materials make use of a mathematical representation of the atomic structure that is both spatially localized around the atoms and invariant under a rigid rotation of the system. While generally successful for the prediction of electronic energies, these models carry two main limitations: i) on the one hand, the local nature of the atomistic representation implies neglecting any physical phenomenon that occurs beyond the cutoff distance selected  to define the spatial extent of the atomic environment, ii) on the other hand, many physical contexts require to deal with properties that are not necessarily invariant, but rather covariant, under a three-dimensional rotation of the system. I will first introduce a class of long-range structural features that can be used to bypass the nearsightedness of local representations by incorporating a Coulomb-like non-local description within finite-size atom-centered environments. I will show that these kind of features not only allow us to rigorously capture the electrostatic energy of a system, but also enable the prediction of a broad range of long-range effects, going from electronic polarization to dispersion interactions. I will then introduce a class of covariant representations of the atomic structure that present the same symmetry as spherical harmonics. In particular, I will show how to use these structural features to efficiently regress the irreducible decomposition of tensors of arbitrary rank, such as dipole moments and  polarizabilities, paving the way to the inexpensive simulation of the spectroscopic signature of a material. Finally, I will show how the very same class of symmetry-adapted representations can be used to tackle the more challenging problem of learning the electronic charge density of a system, as well as any other scalar field, in a highly transferable and linear-scaling fashion."""
\section{Representing Atomic Densities}

We can represent a molecule in a density representation by the equation (using physics style notation).

\begin{align}
    \langle r | \mathcal{A} \rangle &= \sum_i g(r-r_i)|\alpha_i \rangle
\end{align}

Here $|\alpha_i \rangle$ is a feature vector that describes the local atomic environment around atom $i$, and $g$ is some interaction term which describes the effect of atom $i$ on position $r$ (typically some Gaussian). In general, feature vectors $|\alpha_i \rangle$ are computed using local neighbors within some distance cutoff. To expand this cutoff to allow for long range interactions, we need to allow for a nonlocal interaction term. We perform the update
\begin{align}
g(r-r_i) \mapsto \int dr' \frac{g(r'-r_i)}{|r' - r|^p}
\end{align}
This update replaces the local feature vector with an integral over all space. This integral, especially for low values of $p$ enables long range communication between atoms. After this transformation, the atom density potential representation is given by
\begin{align}
    \langle r | \mathcal{V}^p \rangle &= \sum_i |\alpha_i \rangle \int dr' \frac{g(r'-r_i)}{|r'-r|^p
    }
\end{align}

This basic representation is not invariant to transformation groups. To provide invariance over group $G$, we perform the integral

\begin{align}
\int dG \langle r | G | \mathcal{V}^p \rangle
\end{align}

By symmetrizing over the rotation, translation and inversion groups, we can obtain a representation equivariant to $E(3)$.

\section{Symmetry adapted tensor products}

Basis set methods represent the electronic density $\rho(r)$ as a sum of atom-centered basis functions

\begin{align}
    \rho(r) &= \sum_i \rho_i(r) \\
    &= \sum_{i, k} c^i_k \phi^i_k(r) \\
    &= \sum_{i, k} c^i_k \phi_k(r-r_k)
\end{align}
Here the $\phi_k$ are a basis set of atom-centered density functions. Each basis function is factorized as

\begin{align}
    \phi_k(r) &= \phi_{nlm} \\
    &= R_n(|r|)Y^l_m(\hat{r}) 
\end{align}
where the $R_n$ are radial basis functions and $Y^l_m$ are spherical harmonics. 

%provide rotational averages.?

%2-3-4 body components. SOAP is the dual formulation of the rotational tensorial effects. Basically a lot of existing energy functions really break down past the local energy domain.Some ideas, fixed-point charges, vdW corrections, or ML learning of  partial charges, dipoles, multipoles, electronegatativity.

%\section{Long Distance Equivariant Representations}
\section{Training A LODE Model}

As with ACE, the LODE representation computes plain features that are non-tunable, so the method can be fitted directly with linear regression.

%This work introduces LODES, long distance equivariants.

%Implement a potential harmonics function. We make linear models for long range energies. You can do some very interesting things like dipole-quadrupole interactions. You can even do things like model mutual polizarization between a water molecule and a metallic slab of lithium.

%\textcolor{red}{TODO: 
%Basic idea: construct nonlocal representations of feature vectors that are equivariant in O(3).
%Covariant regression: 
%Irreducible spherical decompositions.
%Rotational covariance comes from angular momentum composition.
%SALTED packed.
%github Tensoap.
%Discussion, the impact of learned functionals could be immense.
%}
\end{document}